This is a small CPU study of one synthetic family. (i) Scale: $n=900$, 3 classes, 16-d Gaussian features, one mean degree for the main grid, two layers, hidden size 32; conclusions about where crossovers lie depend on these and need not transfer to larger or real graphs. (ii) Seeds: 5 seeds per cell, so cell-level $p$-values are noisy; the tests are uncorrected for the many cells examined and the verdicts we call supported or refuted at single cells (e.g.\ $h=0.4$, $\mu=1$) could flip with more seeds. (iii) Reimplemented baselines: our H2GCN-style model is a simplified ego/neighbour version without the intermediate-representation concatenation of the published method, our GCN and LP are reimplemented, and nothing was checked against the authors' code. (iv) Tuning: a 8-configuration grid on separate graphs at three cells with one feature strength, shared across all $h$ and $\mu$; per-cell tuning, other depths or hidden sizes could change rankings, especially the MLP/GCN margin at mid homophily. (v) Model: uniform inter-class edge probabilities make heterophily maximally informative about class (neighbours are always another class), which flatters models that can use neighbour information; structured inter-class edges, degree heterogeneity and class imbalance are not covered. (vi) Mechanisms: we did not run experiments isolating why GCN loses to the MLP or why tied weights fail, so those explanations are hypotheses. (vii) Post-hoc choice: the crossover reading used for H4 was clarified after seeing the data. A full study would add real heterophilous benchmarks~\cite{platonov2023critical}, more seeds, other model families~\cite{chien2020adaptive,lim2021large}, and per-cell tuning.
